% page 368 of the facsimile (image p-367 of the scan)
% block 154.9 x 237.7 mm ; left margin 8.0 mm
\pgc{-12.700mm}{0.000mm}{
\l{\cel{2}360.}
\l{}
\l{\sou{masto}. I. (\sou{navig.}) Longa peco de ligno, cilindra e rekta,}
\l{\cel{3}muntita vertikale sur navo, barko, e destinita a supor-}
\l{\cel{3}tar la yardi an qui esas ligita la segli. - II. Ligno-}
\l{\cel{3}peco plu o min longa, muntita vertikale, en gimnastik-}
\l{\cel{3}eyo por exerco pri klimo. - DEFRS.}
\l{}
\l{\sou{mastico}. Nomo di substanci diversa, kompozuro, mi-solida,}
\l{\cel{3}quin onu uzas por stopar trui, fenduri, junturi, e qui}
\l{\cel{3}hardeskas segun ke tempo fluas. - DEFI.}
\l{}
\l{\sou{mastikar}. (\sou{trans.}) Pecigar per la denti, per la movado di}
\l{\cel{3}la maxilo sur la mandibulo, (la nutrivi solida) por igar}
\l{\cel{3}li plu facile glutebla e digestebla. - EFIS.}
\l{}
\l{\sou{mastodonto}. (\sou{paleontol.}) Mamifero fosila, proxima a la}
\l{\cel{3}elefanto. - DEFIRS.}
\l{}
\l{\sou{mastro}.\cel{2}I. La persono qua dominacas personi, kozi. -}
\l{\cel{3}II. La persono qua dominacas en ta o ca domeno.}
\l{\cel{3}- DEFIS.}
\l{}
\l{\sou{mato}. Texuro de palio, de junko tresigita, per qua onu}
\l{\cel{3}kovras la planko-sulo, la muri. - DE.}
\l{}
\l{\sou{matadoro}. (\sou{precipue en Hispania)}. La viro qua, en la tauro-}
\l{\cel{3}kombati, havas kom tasko ocidar la animalo per espado-}
\l{\cel{3}stroko. - DEFIRS.}
\l{}
\l{\sou{matematiko}. La cienco (aritmetiko, geometrio o mekaniko)}
\l{\cel{3}di qua la skopo esas la mezuro di la grandori.-DEFIRS.}
\l{}
\l{\sou{materio}. I.Substanco ye qua kozo konsistas. - II. To quo}
\l{\cel{3}esas la objekto di la agado o facado da la homo. -}
\l{\cel{3}III. (\sou{filoz.}) Substanco korpa. (antonimo di \textquotedbl{}spirito\textquotedbl{}).}
\l{\cel{3}- DEFIRS.}
\l{}
\l{\sou{matida}. I. Qua ne brilas.- II. (\sou{bruiso}): Poka sonora. -}
\l{\cel{3}III. (\sou{koloro}): Qua ne dazlas, nek mem nur brilas. -}
\l{\cel{3}FS.}
\l{}
\l{\sou{matino}. La tempo qua fluas inter jornesko e dimezo. - EFIS.}
\l{}
\l{\sou{matloto}. (\sou{koquarto}).Disho ek plura sorti de fishi, pecigita,}
\l{\cel{3}preparita per vino. - F.}
\l{}
\l{\sou{matro}. Patrino homa.(Equivalanto sentimentala di \textquotedbl{}patrino\textquotedbl{}.)}
\l{}
\l{\sou{matraco}. Longa kuseno, polsterizita, plenigita ye lano, buro,}
\l{\cel{3}krino, e c., quan onu extensas sur lito e sur qua onu su}
\l{\cel{3}kushas. - DEFIR.}
\l{}
\l{\sou{matraso}. Vazo de vitro, di qua la kolo esas longa e streta,}
\l{\cel{3}quan onu uzas en farmacio, kemio. - EFIS.}
}
